\documentclass[final]{siamltex}
\usepackage{graphicx,amsmath,amssymb,latexsym}

\title{A Survey of Applications of the Julia Variation}

\author{Roger W. Barnard, Kent Pearce \thanks{Department of 
Mathematics and Statistics,  Texas Tech University, 
Lubbock, Texas  79409 
({\tt barnard@math.ttu.edu, pearce@math.ttu.edu}) }
\and Carolynn Campbell\thanks {Department of 
Mathematics, Austin Community College, 
Austin, Texas  }}

\begin{document}

\maketitle
\begin{abstract}

 \noindent This is an introductory survey on applications of the Julia variation to problems in geometric 
function theory. A short exposition is given which develops a method for treating extremal problems
over classes $F$ of analytic functions on the unit disk $D$ for which appropriate subsets $F_n$ can be 
constructed so that (i) $F = \overline{\bigcup_n F_n}$ and 
(ii) for each $f \in F_n$ a geometric constraint will hold that $\partial f(D)$ will have at most $n$ ``sides".
Applications of this method which have been made to problems in the literature are reviewed, e.g., Netanyahu's 
problems about the distortion theorems for starlike and convex functions constrained to contain a fixed disk; 
Goodman's problems about omitted values for classes of univalent functions; integral means estimates 
for derivatives of convex functions; maximization problems for functionals on linear fractional transforms 
of convex and starlike functions.
\end{abstract}


\begin{keywords} 
Julia variation, variational methods, geometric function theory 
\end{keywords}

\begin{AMS}
30C
\end{AMS}

\pagestyle{myheadings}
\thispagestyle{plain}
\markboth{R.W. Barnard, C. Campbell, K. Pearce}{A Survey of Applications of the Julia Variation}




This is a survey paper on the applications of a variational method introduced by J. Krzyz in \cite{Krzyz1}.
It is based on Julia's modification of Hadamard's variation of the Green's function.  In the early
1900's, in order to investigate the behavior of Green's function for slight deformations on the boundary 
of its domain, Hadamard developed his variational formulas \cite{Haramard1}.  The validity of the formulas 
relied on the domain having (at least) a continuously differentiable or smooth boundary.  In the 1930's 
Julia, modifying Hadamard's procedure, derived a new variational formula in terms of the Riemann mapping
function for the domain.  This also required the boundary of the domain to be smooth.  This requirement 
proved to be too restrictive to solve most extremal problems since the corresponding extremal domains 
have non-smooth, typically piece-wise smooth, boundaries.  In particular, Schiffer suggested in \cite{Schiffer1} that the requirement of
a smooth boundary was too restrictive for extremal problems for the general class of univalent functions; he 
then proceeded to develop his method of interior variations to circumvent the issue.


The first author was introduced to the Hadamard and Julia variational formulas during his graduate studies 
at the University of Maryland by J. Hummel.  Hummel had used results on the Julia variation, obtained in
\cite{SchSchSp1}, to create in \cite{Hummel1} a general variational method for starlike functions.  He gave a
discussion of the Hadamard and Julia formulas in his lecture notes in \cite{Hummel2}.  We note also that 
Julia's formula was used by M.S. Robertson in \cite{Robertson1} to develop a variational method for analytic
functions on the unit disk with positive real part.


Aside from Schiffer's and Hummel's references, the Hadamard and Julia variational formulas remained generally 
dormant until Krzyz in \cite{Krzyz1} 
applied the Julia variational formula to convex polygons.  A rigorous proof that the formula is valid when 
the boundary contains corners was provided by Barnard and Lewis in \cite{BarnardLewis1}.  It was found later that an 
independent proof had been given earlier by Warshawski in \cite{Warshawski1}.  Consequently, the initial condition
of smoothness on the boundary of the domain could be replaced by a piece-wise smooth condition.  This allowed 
for the method to be applied to a fairly large class of functions; in particular, to any class of functions whose 
image domains could be approximated by domains bounded by a finite number of arcs with explicit geometric conditions.
The extremal domains with their piece-wise smooth boundaries are then among the approximating domains with a small
number of sides. Since the approach of using dense subclasses does not usually give uniqueness of the extremal 
functions other methods are needed to obtain uniqueness.

For completeness, we will give a brief outline of the method.  Let $\Bbb  C$ denote the complex plane and let
\begin{align*}
D_r &= \{ z : |z| < r \}, \hspace{0.5in}D = \{ z : |z| < 1 \} = D_1, \\
C_r &= \{ z : |z| = r \}, \hspace{0.5in}\gamma = \{ z : |z| = 1 \} = C_1, \\
{\cal A} &= \{ f : f \mbox{ is analytic on } D, \ f(z) = z + . . . \ \}, \\
S &= \{ f : f \in {\cal A}, f \mbox{ is 1-to-1 on } D \} \\
S_X &= \{ f : f \in {\cal A}, f(D) \mbox{ is geometrically characterized by } X \}.
\end{align*} 

Two specific subclasses of $S$ to which we will make frequent reference are $S^*$ and $K$, the subclasses of starlike functions 
(with respect to the origin) and convex functions, respectively. Recall that a function $f \in S$ is starlike if and only if 
$${\rm Re} \  \{ \frac{z f'(z)}{f(z)} \} > 0 $$

\noindent for $z \in D$. Also, a function $f \in S$ is convex if and only if 
$${\rm Re} \  \{ 1 + \frac{z f''(z)}{f'(z)} \} > 0 $$

\noindent for $z \in D$.

Let $E$ be a domain in $\Bbb  C$.  Let $r_0$ be the {\em mapping radius} or {\em inner radius} of the domain $E$ at $z_0$.
It is determined by the limit
$$ \lim_{z \to z_0} g(z,z_0) + \log |z - z_0| = \log r_0$$

\noindent where $g(z,z_0)$ is the Green's function of $E$ at $z_0$.  Alternately, the mapping radius is given as
the first order coefficient of the univalent mapping function $f : D \rightarrow E$ such that $f(0) = z_0$, i.e., let 
$f(z) = a_0 + a_1 z + a_2 z^2 + . . . $, then $r_0 = |a_1|$.
Both the mapping radius and the Green's function depend on the
domain (and $z_0$), and as was shown in \cite{Hayman1}, the mapping radius is monotonically, set theoretically and 
continuously dependent on the size of the domain.

Let $X_0$ be a specific geometric characterization.  We will suppose that there exist compact subclasses
$$S_n = \{ f \in S_{X_0} : \partial f(D) \mbox{ has at most } n \mbox{ smooth ``sides"} \}$$

\noindent such that for each $f \in S_n, \partial f(D)$ is a bounded Jordan curve and 
$$\overline{\ \bigcup\limits_n S_n} = S_{X_0}.$$

\noindent The definition of ``sides" will depend on the geometric characterization of $X_0$.

For example, for the class of convex functions $K$ the subclass $K_n$ could be all of the functions for which $f(D)$ is a bounded
convex polygon with at most $n$ sides.

Let $X_0$ be a specific geometric characterization.  Suppose $L$ is a continuous linear functional, $L : {\cal A} \rightarrow {\Bbb C}$, and for the 
class $S_{X_0}$ we want to find $\sup_{f \in S_{X_0}} {\rm Re} \  L(f)$.  Standard normal family arguments give the
existence of an extremal function whenever $S_{X_0}$ is a normal family.  

The following proposition will now be proven.

\vspace{0.1in}

{\bf Proposition 1.} {\em Let $z = re^{i \theta}, |z| = r < 1$, and let $T(\zeta) = L \left ( z f'(z) \frac{1+\zeta z}{1-\zeta z} \right )$, where
$\zeta = e^{i \phi}$.  If ${\rm Re } \  T(\zeta) = c$ (for any constant $c$) has at most $m$ solutions on $|\zeta| = 1$, 
then for $n \geq m$
$$\sup \limits_{f \in S_{X_0}} {\rm Re} \  L(f) = \sup \limits_n \max \limits_{f \in S_n} {\rm Re} \  L(f) = \max \limits_{f \in S_m} {\rm Re} \  L(f).$$
}

{\bf Proof.}  The compactness of the subfamilies $S_n$ along with the density of $\bigcup_n S_n$ in $S_{X_0}$ 
insures that, for $n \geq m$, 
$$\sup \limits_{f \in S_{X_0}} {\rm Re} \  L(f) = \sup \limits_{n \geq m} \max \limits_{f \in S_n} {\rm Re} \  L(f).$$

We will show that $\max_{f \in S_n} {\rm Re} \  L(f) = \max_{f \in S_m} {\rm Re} \  L(f)$ for $n \geq m$.

Let $n$ be a fixed integer greater than $m$.  Consider any function $f$ in $S_n$ where $\partial f(D)$ has strictly more than $m$ ``sides".  By the 
assumption that $\partial f(D)$ is a (bounded) Jordan curve (where $f(D)$ is a Jordan domain), we have, by applying Caratheodory's Extension
Theorem \cite{Duren1}, that $f$ has a continuous extension to $\bar{D} = D \bigcup \gamma$.

Let $\Omega = f(D)$ and $\Gamma = f(\gamma) = \partial \Omega$.  Then $\Gamma$ has $m+1$ or more ``sides."  Call these sides, $\Gamma_1, \Gamma_2, . . . $ and call the
preimages $\gamma_1, \gamma_2, . . .$, where $f(\gamma_j) = \Gamma_j$.

We will locally vary $\Gamma$ by moving one ``side" in the following manner (see Figure 1).  Pick a ``side" of $\Gamma$, say $\Gamma_j$.  
For $w \in \Gamma_j, w = f(z)$ for some $z \in \gamma_j$.  Let $n(w) = \frac{zf'(z)}{|zf'(z)|}$ and let $p(w)$ be the distance
along the unit normal from $w$ to a point $w^*$ on a new ``side" $\Gamma_j^*$.  The distance moved is controlled by $\epsilon p(w)$, which is 
positive along the outward pointing normal and negative along the inward pointing normal.  That is, it is positive if the ``side" is moved out
and negative if the ``side" is moved in.  (The variational distances at the corners are shown to be $o(\epsilon)$ in \cite{BarnardLewis1}.)  Form a new polygonal curve
$\Gamma_\epsilon$ containing the ``side" $\Gamma_j^*$.  When the new curve $\Gamma_\epsilon$ can be chosen so that it still has $n$ or more ``sides,"
the variation is allowable.  Details can be found in Barnard \cite{Barnard2}.

$$\includegraphics[scale=0.5]{var-1a-n} \hspace{0.5in} \includegraphics[scale=0.5]{var-1b-n}$$

\begin{center}Figure 1\end{center}

\vspace{0.1in}

We note here that our initial definition of ``sides" will require a geometric condition which allows the ``sides" to be moved
both in and out in the manner described above.  

For a given $z \in D \ (z = re^{i \theta})$, the varied function $f_\epsilon$  which maps $D$ onto the varied domain $\Omega_\epsilon$ for which 
$\partial \Omega_\epsilon = \Gamma_\epsilon$ becomes by the Julia Variational Formula \cite{Barnard2},
$$f_\epsilon(z) = f(z) + \epsilon z f'(z) \int_\Gamma \frac{1+\zeta z}{1 - \zeta z} d\Psi + o(\epsilon)$$

\noindent where $\zeta = e^{i\phi}, o(\epsilon)$ is uniform for $z$ in compact subsets of $D$ and, for $w = f(\zeta)$,
$$d\Psi = \frac{p(w)n(w)}{i[\zeta f'(\zeta)]^2}dw.$$

A change in variable
$$d\Psi(\phi) = \frac{p(w)}{|f'(\zeta)|}d\phi$$

\noindent insures that $d\Psi(\phi)$ is real and gives
$$f_\epsilon(z) = f(z) + \epsilon z f'(z) \int_\gamma \frac{1+\zeta z}{1 - \zeta z} d\Psi(\phi) + o(\epsilon).$$

\noindent We note for future reference that the change in mapping radius from $f(D)$ to $f_\epsilon(D)$ is given by
$$\epsilon \int_\gamma d\Psi(\phi) + o(\epsilon).$$

\noindent If the varied side is $\Gamma_j$, the integral need only be taken over $\gamma_j$ since the value of 
$d\Psi(\phi)$ is defined to be zero over the rest of $\gamma$ by construction.

If two ``sides" of $\Gamma$ are locally varied, say $\Gamma_k$ was moved out and $\Gamma_l$ was moved in, the varied function
becomes
$$f_\epsilon(z) = f(z) + \epsilon \left [ z f'(z) \int_{\gamma_k} \frac{1+\zeta z}{1 - \zeta z} d\Psi
- z f'(z) \int_{\gamma_l} \frac{1+\zeta z}{1 - \zeta z} d\Psi \right ] + o(\epsilon). $$


Using the linearity of $L$, we have
\begin{align*}
{\rm Re} \  L(f_\epsilon) = {\rm Re} \  L(f) &+ \epsilon \int_{\gamma_k} {\rm Re} \  L \left ( z f'(z) \frac{1+\zeta z}{1 - \zeta z} \right ) d\Psi \\
&- \epsilon \int_{\gamma_l} {\rm Re} \  L \left ( z f'(z) \frac{1+\zeta z}{1 - \zeta z} d\Psi \right ) + o(\epsilon).
\end{align*}

\noindent We have then
\begin{align*}
\frac{\partial}{\partial \epsilon}  {\rm Re} \  L(f_\epsilon)|_{\epsilon=0}  &= \int_{\gamma_k} {\rm Re} \  L \left ( z f'(z) \frac{1+\zeta z}{1 - \zeta z} \right ) d\Psi \\
&- \int_{\gamma_l} {\rm Re} \  L \left ( z f'(z) \frac{1+\zeta z}{1 - \zeta z} d\Psi \right ).
\end{align*}

Let $T(\zeta) = L(z f'(z) \frac{1+\zeta z}{1- \zeta z})$.  Then,
\begin{equation}
\frac{\partial}{\partial \epsilon}  {\rm Re} \  L(f_\epsilon)|_{\epsilon=0}  = \int_{\gamma_k} {\rm Re} \  T(\zeta) d\Psi
- \int_{\gamma_l} {\rm Re} \  T(\zeta) d\Psi. \label{eq1}
\end{equation}

The function $T(\zeta)$ maps $|\zeta| = 1$ onto a curve $\beta$.  On this curve $\beta$, there are $m+1$ arcs resulting from 
$T(\gamma_1), T(\gamma_2), . . . , T(\gamma_{m+1})$ (see Figure 2).  The Mean Value Theorem for integrals allows us to conclude 
that on each arc $\gamma_j$ there exists a $\zeta_j$ such that 
\begin{equation}
\int_{\gamma_j} {\rm Re} \  T(\zeta) d\Psi(\phi) = {\rm Re} \  T(\zeta_j) \int_{\gamma_j} d\Psi(\phi).  \label{eq2}
\end{equation}

$$\includegraphics[scale=0.5]{var-2-n}$$

\begin{center}Figure 2\end{center}

\vspace{0.1in}

This implies that for $\gamma_k$ and $\gamma_l$ there are $\zeta_k$ and $\zeta_l$ satisfying (\ref{eq2}) which makes (\ref{eq1})
become
$$\frac{\partial}{\partial \epsilon}  {\rm Re} \  L(f_\epsilon)|_{\epsilon=0}  = {\rm Re} \  T(\zeta_k) \int_{\gamma_k}  d\Psi
- {\rm Re} \  T(\zeta_l) \int_{\gamma_l}  d\Psi. $$

By hypothesis , ${\rm Re} \  T(\zeta) = c$ (for any constant $c$) has at most $m$ solutions.  Thus, the most such points, $T(\zeta_j)$, 
any vertical line (${\rm Re} \ T(\zeta) = c$) could intersect is $m$.  Since $\partial f(D)$ has strictly more than $m$ sides, there must exist two points
$T(\zeta_k)$ and $T(\zeta_l)$ in $\beta$ such that the real part of one is greater than the real part of the other, say
$${\rm Re} \  T(\zeta_k) > {\rm Re} \  T(\zeta_l).$$

\noindent Thus, if the two sides that are initially varied are $\Gamma_k$ and $\Gamma_l$, then the partial in (\ref{eq1}) can
be chosen so that 
$$\frac{\partial}{\partial \epsilon} {\rm Re} \  L(f_\epsilon)|_{\epsilon=0} > {\rm Re} \  T(\zeta_l)
\left [ \int_{\gamma_k} d\Psi(\phi) - \int_{\gamma_l} d\Psi(\phi) \right ].$$

As noted earlier, $\int_{\gamma_k} d\Psi(\phi) - \int_{\gamma_l} d\Psi(\phi)$ represents the change in the mapping radius.  From
continuity $\Gamma$ can be varied so that the resulting function $f_\epsilon$ is again in $S_n$, i.e., if $\Psi(\phi)$ is chosen on 
$\gamma_k$ and $\gamma_l$ so that the resulting mapping radius of $f_\epsilon$ is 1, then $\int_{\gamma_k} d\Psi(\phi) - \int_{\gamma_l} d\Psi(\phi) = 0$.
Hence, 
$$\frac{\partial}{\partial \epsilon} {\rm Re} \  L(f_\epsilon)|_{\epsilon=0} > [{\rm Re} \  T(\zeta_l)](0) = 0.$$

Recall that
\begin{align*}
\frac{\partial}{\partial \epsilon}  {\rm Re} \  L(f_\epsilon)|_{\epsilon=0}  &= \int_{\gamma_k} {\rm Re} \  L \left ( z f'(z) \frac{1+\zeta z}{1 - \zeta z} \right ) d\Psi \\
&- \int_{\gamma_l} {\rm Re} \  L \left ( z f'(z) \frac{1+\zeta z}{1 - \zeta z} d\Psi \right ).
\end{align*}

\noindent We have
$${\rm Re} \  L(f_\epsilon) = {\rm Re} \  L(f) + \epsilon \left [ \frac{\partial}{\partial \epsilon} 
{\rm Re} \  L(f_\epsilon)|_{\epsilon=0} \right ] + o(\epsilon).$$

\noindent Thus, for any function $f \in S_n$ such that $\partial f(D)$ has strictly more than $m$ ``sides,"
each of which can be moved both in and out, a function can be found which increases the real part of the functional $L$.
Therefore, we have that the function $f \in S_n$ which maximizes the real part of the functional $L$ is actually in $S_m$.
Thus, for $n \geq m$, 
$$\max \limits_{f \in S_n} {\rm Re} \  L(f) = \max \limits_{f \in S_m} {\rm Re} \  L(f).$$

This proves Proposition 1.

\vspace{0.1in}

Remark.  In \cite{Barnard6} it is shown that by combining this method with Loewner theory on slit domains 
that this method may be applied to domains where $\partial f(D)$ is not a Jordan curve (where $f(D)$ is a 
slit domain) with slight variations.  Hence, this variational method involving Proposition 1 may be used 
to solve extremal problems 
over a much broader class of functions.

It is this basic idea that the authors have been able to apply to a variety of extremal problems.  The first 
arose in Barnard's dissertation and answered a question which arose out of Netanyahu's work in \cite{Netanyahu1}.  
Netanyahu had asked to determine the extremal functions for the standard distortion theorems for the class 
of starlike or convex functions whose image domains contain a fixed disk centered at the origin.  Netanyahu 
had done this \cite{Netanyahu1} for the class $NS(d)$ of univalent functions in $S$ whose image domains contain 
a fixed disk of radius $d$ centered at the origin, but his methods did not appear to work for the corresponding 
classes $NS^*(d)$ or $NK(d)$.  The extremal functions for the classes $NS^*(d)$ and $NK(d)$ were found by 
Barnard in \cite{Barnard2}, however, a more general result containing resolutions for Netanyahu's questions 
as special cases was given in \cite{BarnardLewis1} by Barnard and Lewis.

To obtain the more general result, we first answered M. Reade's question as to determining a geometric 
characterization for the Mocanu class, $MS^*(\alpha)$, of $\alpha$-starlike functions.  $MS^*(\alpha)$ 
was a continuously parametrized family of functions, $0 \leq \alpha \leq \infty$, defined analytically, 
which contained the class $S^* = MS^*(0)$ and the class $K = MS^*(1)$.  In particular, Mocanu had defined 
in \cite{Mocanu1} the class $MS^*(\alpha)$ as those functions in ${\cal A}$ for which 
$${\rm Re} \  \left \{ (1-\alpha) \left (\frac{z f'(z)}{f(z)} \right ) + \alpha \left ( 1 + \frac{z f''(z)}{f'(z)} \right ) \right \} > 0.$$

\noindent We determined in \cite{BarnardLewis1} that by defining an $\alpha$-arc for $\alpha > 0$, 
loosely speaking, as a translate of an arc of the $\alpha^{\rm th}$ power of the line $x = 1$ (see \cite{BarnardLewis1} for details), 
then a domain $f(D)$ is the image of an $\alpha$-starlike function $f$ if and only if any two 
points in $f(D)$ can be connected by an $\alpha$-arc lying 
in $f(D)$.  (Note $\alpha = 1$ gives $1$-arcs as straight line segments [hence, convexity], while the limiting case of 
$\alpha=0$ gives starlikeness with respect to the origin.)  By using as a dense subclass those domains 
bounded by a finite number of $\alpha$-arcs as ``sides", we proved the following result.

For a given $d, 0 \leq d < 1$ and $1 < \rho \leq \infty$, let $MS^*(\alpha,d,\rho)$ be those functions in $MS^*(\alpha)$ for 
which $d \leq |f(z)/z| \leq \rho$.  Let $F = F(\cdot,\alpha, d,\rho)$ be the function in $MS^*(\alpha,d,M)$ whose image is the symmetric domain bounded by an arc $\gamma_d \subset C_d$, which is symmetric about $-d$, an arc $\gamma_\rho \subset C_\rho$, which is symmetric about $\rho$ and two $\alpha$-arcs connecting the endpoints of $\gamma_d$ and $\gamma_\rho$. We proved the following main theorem and corollary.

\vspace{0.1in}

{\bf Theorem 1.} {\em Let $\alpha, d$ and $\rho$ be fixed nonnegative numbers satisfying $0 \leq \alpha \leq \infty$, $0 \leq d < 1$ and $1 < \rho \leq \infty$. Then, 

(A) The function $g(z) = \log [F(z)/z], z \in D$ is univalent and convex in the direction of the imaginary axis.

(B) If $f \in MS^*(\alpha,d,\rho)$, then $\log [f(z)/z]$ is subordinate to $g$.}

\vspace{0.1in}

{\bf Corollary 1.} {\em Let $\alpha, d$ and $\rho$ be as in Theorem 1.  Let $\Phi$ be a given nonconstant entire
function.  If $f \in MS^*(\alpha,d,\rho)$, then

(A) For given $z \in D \backslash \{0\}$ 
$${\rm Re} \  \left \{ \Phi \left [ \log \frac{f(z)}{z} \right ] \right \} \leq 
\max \limits_{0 < \theta \leq 2 \pi} {\rm Re} \  \left \{ \Phi \left [ \log \frac{F(z)}{z} \right ] \right \},$$

(B) For given $r, 0 < r < 1$, and $\lambda > 0$,
$$\int_{-\pi}^\pi |f(re^{i \theta})|^\lambda d \theta \leq \int_{-\pi}^\pi |F(re^{i \theta})|^\lambda d \theta,$$

(C) For a given positive integer $N \geq 2$,
$$\sum_{k=2}^N |a_k|^2 \leq \sum_{k=2}^N |A_k|^2,$$

\vspace{0.1in}

\noindent where $f(z) = z + \sum_{k=2}^\infty a_k z^k$ and $F(z) = z + \sum_{k=2}^\infty A_k z^k, z \in D$. }


Another application of this variational method was to show in \cite{Barnard4} that any $f \in S$ for which $\partial f(D)$
has an analytic slit which contains a point $f(e^{i \theta_1}) = f(e^{i \theta_2})$ where the opposing
normals exist and are of unequal modulus, i.e., $|f'(e^{i \theta_1})| \neq |f'(e^{i \theta_2})|$, can always be varied 
locally on the boundary to produce a domain of strictly larger mapping radius.  For consequences of this result, see 
\cite{Barnard4}.


General applications of the variational method were used in a series of papers on the omitted values problem of Goodman.  The problem of omitted values was first posed by Goodman \cite{Goodman1} in 1949, restated by MacGregor \cite{MacGregor1} in his survey article in 1972, then reposed in a more general setting by Brannan \cite{AndBarBran1} in 1977.  It also appears in Bernardi's survey article \cite{Bernardi1} and has 
appeared in several open problem sets since then, including \cite{Barnard1}, \cite{BrannanClunie1}, \cite{Goodman2} and \cite{Miller1}.


For a function $f \in S$, let $A(f)$ denote the Lebesgue measure of the set $D \backslash f(D)$ and for $0 < r < 1 $ let $L(f,r)$ denote
the Lebesgue measure of the set $\{D \backslash f(D)\} \bigcap C_r$.  Two explicit problems posed by Goodman and by Brannan were to determine
\begin{equation*}
A = \sup_{f \in S} A(f) 
\end{equation*} 

\noindent and
\begin{equation}
L(r) = \sup_{f \in S} L(f,r). \label{eq4}
\end{equation}

Goodman \cite{Goodman1} showed that $0.22\pi < A < 0.50\pi$.  The lower bound which he obtained was generated by a domain
of the type shown in Figure 3.

$$\includegraphics[scale=0.75]{var-3-n}$$

\begin{center}Figure 3\end{center}

\vspace{0.1in}

Later, Goodman and Reich \cite{GoodmanReich1} gave an improved upper bound of $0.38\pi$ for $A$.  Lewis had shown in \cite{Lewis1}
by applying some deep results of Alt and Caffarelli \cite{AltCaffarelli1} in partial differential equations for free boundary
problems that the extremal domains had to have piecewise analytic boundaries and that our methods developed in 
\cite{BarnardLewis1} for these omitted value problems could be used to give a geometric description for the boundaries of the extremal domains as follows. There is 
an $f_0$ in $S$ with $A = A(f_0)$ such that $f_0(D)$ is circularly symmetric with respect to the positive real axis, i.e., it has the
property that for $0 < r < 1$,
$$\frac{\partial}{\partial \theta} |f_0(re^{i \theta})| \leq 0 \mbox{ and } 
\frac{\partial}{\partial \theta} |f_0(re^{-i \theta})| \leq 0, \mbox{ for } 0 < \theta < \pi$$

\noindent (cf. Hayman \cite{Hayman1}).  Moreover, the $\partial f_0(D)$ consists of the negative real axis up to
$-1$, and arc $\sigma$ of the unit circle that is symmetric about $-1$ and an arc $\lambda$ lying in $D$, except for its endpoints.
The arc $\lambda$ is symmetric about the reals, connects to the endpoints of $\sigma$ and has monotonically decreasing modulus 
in the closure of the upper half disc.  These results follow by standard symmetrization methods.  Much deeper methods are needed 
to show (as in \cite{Barnard5} and \cite{Lewis1}) that $f_0$ has a piecewise analytic extension to $\lambda$ with $f_0'$ continuous
on $f_0^{-1}(\lambda)$ and $|f_0' (f_0^{-1}(w))| \equiv c < 1$ for all $w \in \lambda \bigcap \{ D \backslash (-1,1) \}$.  Using these 
properties of $f_0$ it was shown by the authors in \cite{BarnardPearce1} that by ``rounding the corners" of certain gearlike
domains a close approximation for the extremal function could be obtained.  This gives the best known lower bound of
$$0.24\pi < A.$$

The upper bound is conceptually harder since it requires an estimate on the omitted area for each function in $S$.  Indeed, it appears
difficult to use the geometric description of $f_0$ to calculate $A$ directly.  However, an indirect proof was used by Barnard and Lewis
\cite{BarnardLewis2} to obtain the best known upper bound of 
$$A < 0.31\pi.$$

{\bf Open Problem.}  {\em Show that $f_0$ is unique and determine $A$ explicitly.}

\vspace{0.1in}

For the class $S^*$ of functions in $S$ whose images are starlike with respect to the origin, the problem of determining
the corresponding
$$A^* = \sup_{f \in S^*} A(f)$$

\noindent has been completely solved by Lewis in \cite{Lewis1}.  The extremal function $f_1 \in S^*$ defined by 
$$A^* = A(f_1) \approx 0.235 \pi$$

\noindent is unique (up to rotation).  The extremal domain $\partial f(D)$ is a circularly symmetric domain whose boundary $\partial f_1(D)$ has two radial rays projecting into $D$ with their endpoints connected by an arc $\lambda_1$ that is symmetric about the reals and has $|f_1'(\zeta)| \equiv c_1 \mbox{ for all } \zeta \in f_1^{-1}(\lambda_1)$.

The problem of determining $L(r)$ in (\ref{eq4}) was solved by Jenkins in \cite{Jenkins1} where he proved that for a fixed 
$r, 1/4 \leq r < 1$,
$$L(r) = 2\pi \arccos (8 \sqrt{r} - 8r -1).$$

\noindent The extremal domain in this case is the circularly symmetric domain (unique up to rotation) having as its boundary the negative 
reals up to $-r$ and a single arc of $C_r$ symmetric about the point $-r$.

The corresponding problem for starlike functions of determining 
$$L^*(r) = \sup_{f \in S^*} L(f,r)$$ 
was solved by Lewandowski in \cite{Lewandowski1} and by Stankiewicz in \cite{Stankiewicz1}.  
The extremal domain in that case is the circularly symmetric domain (unique up to rotation) having as its 
boundary two radial rays and the single arc of $C_r$ connecting their endpoints.  An explicit
formula for the mapping function in this case was first given by Suffridge in \cite{Suffridge1}.

However, for the class $K$ of functions in $S$ whose images are convex domains the corresponding problems of determining
\begin{equation*}
A_K(r) = \sup_{f \in K} A(f,r) 
\end{equation*}

\noindent and
\begin{equation*}
L_K(r) = \sup_{f \in K} L(f,r), 
\end{equation*}

\noindent where $A(f,r)$ denotes the Lebesgue measure of $D_r \backslash f(D)$, presents some interesting difficulties.
One particular difficulty is that the basic tool of circular symmetrization used in the solution of each of the previous
determinations is no longer available.  The example of starting with the convex domain bounded by a square shows that 
convexity is not alway preserved under circular symmetrization.  However, Steiner symmetrization (cf. Hayman \cite{Hayman1})
can still be used in certain cases such as sectors.  Another difficulty is the introduction of distinctly different extremal
domains for different ranges of $r$.  Since every function in $K$ covers a disk of radius 1/2 (cf. Duren \cite{Duren1})
$r$ needs only to be considered in the interval $(1/2,1)$.  Waniurski has obtained some partial results in \cite{Waniurski1}.
He defined $r_1$ and $r_2$ to be the unique solutions to certain transcendental equations where $r_1 \approx 0.594$ and $r_2 \approx 0.673$.
If $F_{\pi/2}$ is the map of $D$ onto the half plane $\{ w : {\rm Re} \  w > -1/2 \}$ and $F_\alpha$ maps $D$ onto the sector
$$\left \{ w : \left | arg \left ( w + \frac{\pi}{4\alpha} \right ) \right | < \alpha \right \}$$

\noindent whose vertex, $v = - \pi/4\alpha$, is located inside $D$, then
\begin{align*}
A_K(r) &= A(F_{\pi/2},r) \mbox{ for } 1/2 < r < r_1, \\
L_K(r) &= L(F_{\pi/2},r) \mbox{ for } 1/2 < r < r_1,
\end{align*} 

\noindent and
$$L_K(r) = L(F_\alpha,r) \mbox { for } r_1 < r < r_2.$$

Barnard had announced in his survey talk on open problems in complex analysis at the 1985 {\em Symposium on the Occasion of the 
Proof of the Bieberbach Conjecture} the following conjecture.

\vspace{0.1in}

{\bf Conjecture 1.} {\em The extremal domains for determining $A_K(r)$ and $L_K(r)$ will be half planes, symmetric sectors and domains bounded by single arcs of $C_r$ along with tangent lines to the endpoints of these arcs; the different domains will depend on different ranges of $r$ in $(1/2,1)$.}

\vspace{0.1in}


This conjecture was also made later, independently, by Waniurski at the end of his paper in \cite{Waniurski1} in 1987.

In trying to determine the extremal domains for $A_K(r)$ and $L_K(r)$ via our variational method developed in \cite{BarnardLewis1} the
following problem was investigated.  For $F \subset S$ and $1/4 \leq d \leq 1$ let $F_d = \{ f \in F : \min_{z \in D} |f(z)/z| = d \}$.  Then,
determine the sharp constant $A = A(F_d)$ such that for any $f \in F_d$
\begin{equation}
I_{-1}(f') = \lim_{r \rightarrow 1} \frac{1}{2\pi} \int_{-\pi}^{\pi} \left | \frac{1}{f'(re^{i \theta})} \right | d \theta  \leq \frac{A}{d}. \label{eq7}
\end{equation}

For $0 \leq \alpha < 1$ let $S^*(\alpha)$ denote the subclass of $S$ of starlike functions of order $\alpha$, i.e., a function $f \in S^*(\alpha)$
if and only if $f$ satisfies the condition ${\rm Re} \  z f'(z)/f(z) > \alpha$ for $z \in D$.  It is well known that $K \subset S^*(1/2)$.
It follows fairly easily from subordination theory that $A(S^*(1/2)) \leq 4/\pi$.  Furthermore, this estimate is sharp for $S^*(1/2)$ since the
functions $f_n(z) = z/(1-z^n)^{1/n}$ belong to $S^*(1/2)$ for each $n > 0$.  However, this estimate is not sharp for the class $K$ of convex functions 
which is a proper subset of $S^*(1/2)$.  Considerable numerical evidence suggested to the authors to make the following
conjecture.

\vspace{0.1in}

{\bf Conjecture 2.} {\em For each $d, 1/2 \leq d \leq 1, A = A({K}_d) = 1$ in (\ref{eq7}) with equality holding
for all domains which are bounded by regular polygons centered at the origin.}

\vspace{0.1in}

This conjecture was announced in March 1985 at the {\em Symposium on the Occasion of the 
Proof of the Bieberbach Conjecture} at Purdue University.  It also appeared as Conjecture 8 in the first author's ``Open Problems and Conjectures in 
Complex Analysis" in \cite{Barnard1}.  
It was thought, by many function theorists, that the conjecture would be
easily settled, given the vast literature on convex functions and the
large research base for determining integral mean estimates, see
\cite{Duren1}.

An initial difficulty was the non-applicability of Baernstein's
circular symmetrization methods, since convexity, unlike univalence
and starlikeness, is not preserved under circular symmetrization as noted above.
Although Steiner symmetrization does preserve convexity, see \cite{Valentine1}, it
did not appear to be helpful for the problem and, indeed, we found
that extremal domains need possess no standard symmetry.

A confusing issue, which also arises, is that the integral means of
the standard approximating functions $f_n$ in $K$ defined by
$$f'_n(z)=\prod^n_{k=1}(1-ze^{i\theta_k})^{-2\alpha_k},0<\alpha_k\leq
1,\sum^n_{k=1}\alpha_k=1$$ 

\noindent decrease when the arbitrarily distributed $\theta_k$ are replaced by
uniformly distributed $t_k=k\pi/n$, as was shown in \cite{Zheng1}.  The
conjecture suggests that multiplication by the minimum modulus $d$
must overcome this decrease.

We make the following definition.

\vspace{0.1in}

{\bf Definition.}  Let $\Gamma$ be a curve in $\Bbb C$ such that the
left- and right-hand tangents to the curve $\Gamma$ exist at each
point on $\Gamma$.  The curve $\Gamma$ will be said to {\em
circumscribe a circle} $C$ if the left- and right-hand tangents to
the curve $\Gamma$ at each point on $\Gamma$ lie on tangent lines to
the circle $C$.

\vspace{0.1in}

We will employ the following notation.

\vspace{0.1in}

{\bf Notation.} Let $f\in S$ and suppose that $\lambda$ is a subarc of
$\gamma$ on which $f$ is smooth.  For $z=e^{i\theta}\in\lambda$
let $d_{\theta} = \  <f(z)$, $\frac{zf'(z)}{|zf'(z)|}>$, i.e, $d_\theta$
is the directed length of the projection of $f(z)$ onto the outward
unit normal to the  $\partial f(D)$ at $f(z)$.

\vspace{0.1in}

In 1993 in \cite{BarnardPearce2}, we proved the following theorem which verified 
Conjecture 2.

\vspace{0.1in}

{\bf Theorem 2.} {\em Let $f\in K$, $d={\min\limits_\theta}\ 
|f(e^{i\theta})|$ and $d^*={\sup\limits_\theta} \  d_\theta$.
Then,
\begin{equation*}
\frac{1}{d^\ast}\leq\frac{1}{2\pi}\int^{2\pi}_{0}\frac{1}
{|f'(e^{i\theta})|}d\theta\leq\frac{1}{d} 
\end{equation*}
with equality holding if the  $\partial f(D)$ circumscribes $C_d$. }

\vspace{0.1in}

Our original proof which was quite lengthy used as a major step the following idea which has independent interest.
For any domain $\Lambda$ which contains the origin, let the notation $m.r.(\Lambda)$ denote the inner radius 
of $\Lambda$ at the origin.  Given a $d$ and a convex domain $\Omega = f(D)$ where $f \in {K}_d$, one can always construct
two varied domains $\Omega^*$ and $\Omega^{**}$ which satisfy the containment relationship
\begin{equation}
\Omega^* \subseteq \Omega^{**}, \label{eq9}
\end{equation}

\noindent as follows.

The domain $\Omega^{**}$ is constructed by multiplication by $1 + \epsilon$ for $\epsilon$ sufficiently small
and positive.  This produces a radial enlargement $(1 + \epsilon)\Omega = \Omega^{**}$ of $\Omega$.  The domain $\Omega^*$
is constructed by moving each smooth boundary point $w = f(e^{i \theta})$ a distance $\epsilon d$ in the direction of the
outward pointing normal.  An adjustment is made at the non-smooth boundary points to produce the domain $\Omega^*$ so that 
boundary of $\Omega^*$ has the same geometric restrictions as the boundary of $\Omega$, e.g., same number of sides.  Subordination
then shows that (\ref{eq9}) implies that  $m.r.(\Omega^*) \leq m.r.(\Omega^{**})$.  The variational formulas can then be applied
to yield that
$$ \Delta m.r.(\Omega^*,\Omega) = \frac{\epsilon d}{2\pi} \int_{-\pi}^{\pi} \left | \frac{1}{f'(e^{i \theta})} \right | d \theta  + o(\epsilon) \leq
\epsilon = \Delta m.r.(\Omega^{**},\Omega) $$

\noindent which give the result.  Similar arguments are used for $d^*$.

We obtained, arising out of the proof, the rather unexpected sufficient
condition for equality to occur in (2) for the classes ${K}_d$.
However, because the proof used a scheme to approximate convex
functions by polygonally convex functions, we did not obtain a
necessary condition for equality.

We devised in \cite{BarnardPearce3} a new, simpler proof for the conjecture which extends
Theorem 2.  The proof relaxes the convexity requirement and
validates the necessity of the sufficient condition.

\vspace{0.1in}

{\bf Theorem 3.} {\em Let $f\in S^*(\alpha)$ for some $0 \leq \alpha < 1$.
Suppose $f$ is smooth on $X\subset\partial D$ where $X$ is a
countable union of pairwise disjoint subarcs of $\gamma$ such
that the complement of $X$ in $\gamma$ has measure zero.  Let
$d_*={\inf\limits_{\theta\in X}} \  d_\theta$, $d^*={\sup\limits_{\theta\in X}} \  d_\theta$.  Then,
$$
\frac{1}{d^*}\leq\frac{1}{2\pi}\int^{2\pi}_0
\frac{1}{|f'(e^{i\theta})|}d\theta\leq\frac{1}{d_*}
$$
with equality holding if and only if $\partial f(D)$ circumscribes
$C_{d_*}$. }

\vspace{0.1in}

Another application of the method originated from a question of Clunie and Sheil-Small.  If $f \in S$ and 
$w \not\in f(D)$, then the function
\begin{equation}
\hat{f} = f / (1-f/w) \label{eq10}
\end{equation}

\noindent belongs again to $S$.  The transformation $f \rightarrow \hat{f}$ is an important tool in the study
of geometric function theory.  If $F$ is a subset of $S$, let 
$$\hat{F} = \{ \hat{f} : f \in F, w \in {\Bbb C}^* \backslash f(D) \}.$$

\noindent Here, ${\Bbb C}^* = {\Bbb C} \bigcup \{ \infty \}$.  Since we admit $w = \infty$, it is clear that $F \subset
\hat{F} \subset S$.


If $F$ is compact in the topology of local uniform convergence, then so is $\hat{F}$.  If $F$ is rotationally invariant, that is,
if $f_\alpha(z) = e^{-i \alpha}f(e^{i \alpha}z)$ belongs to $F$ whenever $f$ does, then $\hat{F}$ is also rotationally invariant.  It is 
an interesting question to ask which properties of $F$ are inherited by $\hat{F}$.  Since $\hat{S} = S$, this question is
trivial for $S$.

In \cite{BarnardSchober1}, \cite{BarnardSchober2} and \cite{BarnardSchober3} Barnard and Schober considered transforms of the class $K$ of convex mappings (and the class $S^*$ of starlike mappings).  Simple examples 
show that $\hat{K}$ is strictly larger than $K$.  Since the coefficients of functions in $K$ are uniformly bounded 
[by one (in modulus)], Clunie and Sheil-Small had asked whether the coefficients of functions in $\hat{K}$ have a uniform
bound.  The affirmative solution of this problem was given by Hall \cite{Hall1}.

\vspace{0.1in}

{\bf Open Question.}  {\em Find the best uniform coefficient bound as well as the individual coefficient bounds for functions
in $\hat{K}$.}

\vspace{0.1in}

In \cite{BarnardSchober1} the variational procedure developed in \cite{BarnardLewis1} is applied to a class of extremal problems for 
$\hat{K}$.  If $\lambda : \hat{K} \rightarrow {\Bbb R}$ is a continuous functional that satisfies certain admissibility
criteria, it was shown that the problem
$$\max_{\hat{K}} \lambda$$ 

\noindent has a relatively elementary extremal function $\hat{f}_1$.  More specifically, it was shown that $\hat{f}_1$ is either
a half-plane mapping $\hat{f}_1(z) = z/(1 - e^{i \alpha}z)$ or is generated through (\ref{eq10}) by a parallel strip mapping $f_1 \in K$.

The class of functionals considered in \cite{BarnardSchober1} contain the second-coefficient functional $\lambda(\hat{f}) = {\rm Re} \  a_2$ and the functionals $\lambda(\hat{f}) = {\rm Re} \  \Phi (\log \hat{f}(z)/z)$ where $\Phi$ is entire and $z$ is fixed.  The latter functionals include the problems of maximum and minimum modulus $( \Phi(w) = \pm w)$.  In general, the extremal strip domains $f_1(D)$ need not be symmetric
about the origin.  This adds a nontrivial and interesting character to the problem.  

A sharp estimate for the second coefficient of functions in $\hat{K}$ is given explicitly in the following result.  Surprisingly,
the answer is not an obvious one.

\vspace{0.1in}

{\bf Theorem 4.} {\em If $\hat{f}(z) = z + a_2 z^2 + . . . $ belongs to $\hat{K}$, then
$$|a_2| \leq \frac{2}{x_0} \sin x_0 - \cos x_0 \approx 1.3270$$

\noindent where $x_0 \approx 2.0816$ is the unique solution of the equation
$$\cot x = \frac{1}{x} - \frac{1}{2}x$$

\noindent in the interval $(0, \pi)$.  Equality occurs for the functions $e^{-i \alpha}\hat{f}_1(e^{i \alpha}z), \alpha \in {\Bbb R}$, where
$\hat{f}_1(z) = f_1(z)/[1-f_1(z)/f_1(1)]$ and $f_1$ is the vertical strip mapping defined by
\begin{equation}
f_1(z) = \frac{1}{2i \sin x_0} \log \frac{1 + e^{i x_0} z}{1 - e^{i x_0} z}. \label{eq11}
\end{equation}}

We make the following conjecture.

\vspace{0.1in}

{\bf Conjecture 3.} {\em The extremal functions for maximizing $|a_n|$ over $\hat{K}$ are the vertical strip mappings defined 
by (\ref{eq11}) where a different $x_0$ is needed for each $n$. }

\vspace{0.1in}

In \cite{BarnardSchober2} the Koebe disk, radius of convexity and sharp estimates for the coefficient functional$|t a_3 + a_2^2|$ were found for
functions in the class $\hat{K}$.  Also, in \cite{Ali1} R.M. Ali found sharp upper and lower bounds for $|f(z)|$ for $\hat{f} \in \hat{K}$.


A key lemma in the proofs of the above was the determination of the geometric characterization of the ``sides" to be varied in the 
images of the functions in a dense subclass of the transformed convex domains.  The variations of the functions in ${K}_n$ whose 
image domains are convex polygons with at most $n$ sides did not reduce the problem for $\hat{K}$ to examining domains with a sufficiently
small enough number of sides.  We observed that functions $\hat{f}$ in $\hat{{K}_n}$ map $D$ onto curvilinear polygons with at most $n$
sides and with interior angles at $\pi$.  Furthermore, if $\hat{f} = f / (1-f/w)$, then the sides of the  $\partial \hat{f}(D)$
all lie on circles or lines through the point $w_1 = -w$ (see Figure 4).  In fact, these two properties characterize functions in
$\widehat{{K}_n}$.  That is, if $g \in S$ and the  $\partial g(D)$ is a curvilinear $n$-gon with interior angles at most $\pi$ and if the
sides of the  $\partial g(D)$ all lie on circles or lines through a point $-w \not\in g(D)$, then $f = g/(1-g/w)$ belongs
to ${K}_n$ and so $\hat{f} = g$ belongs to $\widehat{{K}_n}$. 

$$\includegraphics[scale=1.0]{var-4-n}$$

\begin{center}Figure 4\end{center} 

\vspace{0.1in}

These observations thus produced the domains whose boundaries were characterized in a geometric fashion so that our methods
could be applied.

In a more recent application of the Julia variational method J. MA \cite{Ma1} has proved, in his dissertation
under Ruscheweyh, the following results.

Let points $z_j \in D$ with $z_j \neq 0, j = 1, 2, . . . , m$ and positive integers 
$l_j, j = 0, 1, 2, . . . , m$ be given.  Define $N = l_0 + l_1 + l_2 + . . . + l_m$.  For
$g \in {\cal A}$ consider the complex $N$ vector $v(g)$ with components
$$a_k \ (k = 2, 3, . . . , l_0 + 1), \hspace{0.25in} g^{(k)}(z_j) \  (k = 0, 1, . . . , l_j - 1; j = 1, 2, . . . , m).$$
Denote by $V_N(K)$ the set $\{ v(g) : g \in K \}$.  

Let $\Lambda(v)$ be a real-valued function that is differentiable on some open neighborhood of $V_N(K)$.  Define a continuous
functional $\Psi$ on $K$ by 
\begin{align*}
\Psi(g) := \Lambda(v(g)) = \Lambda ( a_2, a_3, . . . , a_{l_0+1}, &g(z_1), g'(z_1), . . . , g^{(l_1 -1)}(z_1), . . .,  \\
 &g(z_m), g'(z_m), . . . , g^{(l_m -1)}(z_m) )
\end{align*}

\noindent A functional of this form will be called a functional of finite degree and the number $N$ is call the degree of the 
functional.  

Let us denote
$$\Lambda_v := \left ( \frac{\partial \Lambda}{\partial w_1}, . . . , \frac{\partial \Lambda}{\partial w_N} \right ),$$

\noindent where
$$\frac{\partial \Lambda}{\partial w_k} = \frac{1}{2} \left ( \frac{\partial \Lambda}{\partial x_k} - i \frac{\partial \Lambda}{\partial y_k} \right ),
\hspace{0.25in} k = 1, 2, 3, . . . , N.$$

The main results are the following theorems.

\vspace{0.1in}

{\bf Theorem 5.} {\em Let $\Psi$ be a functional of degree $N$ on the class $K$ and for $f_0 \in K$ let 
$$\Psi(f_0) = \max (\min) \{\Psi(g) : g \in K \}.$$

\noindent Suppose that $\Lambda_v(v(f_0)) \neq 0$.  Then,
\begin{equation}
f_0 (z) = \int_0^z \frac{d \eta}{ \prod\limits_{k=1}^{N} (1 - z_k \eta)^{2 \lambda_k} }, \label{eq12}
\end{equation}

\noindent where $|x_k| = 1, \lambda_k >= 0, k = 1, . . . , N$ and $\sum_{k=1}^N \lambda_k = 1$. }

\vspace{0.1in}

{\bf Theorem 6.} {\em Let $H(w_1, w_2, . . . , w_N )$ be an analytic function on a neighborhood of the set $V_N(K)$.
If for $f_0 \in K$,
$${\rm Re}\  \{H(v(f_0)) \} = \max (\min) \{\Psi(g) : g \in K \}$$

\noindent then $f_0$ has the form (\ref{eq12}).}

\vspace{0.1in}

We note that extremal functions with the form (\ref{eq12}) map $D$ onto convex polygons with at most $N$ sides.





\newpage


\begin{thebibliography}{99}

\bibitem{Ali1} Ali, R. M. ``A distortion theorem for the class of M\"{o}bius 
transformations of convex mappings.'' {\em Rocky Mountain J. Math.} {\bf 19}
(1989) pp. 1083-1094.

\bibitem{AltCaffarelli1} Alt, H., and Caffarelli, L.  ``Existence and
regularity for a minimum problem with free boundary.''  {\em J. Reine
Angew. Math.} {\bf 325} (1981) pp. 105-144.

\bibitem{AndBarBran1} Anderson, J.M., Barth, K.E., and Brannan, D.A.
``Research problems in complex analysis.''  {\em Bull. London Math
Soc.}  {\bf 9} (1977) pp. 129-162.

\bibitem{Barnard4} Barnard, R.W. ``Criteria for local variations for slit 
mappings." {\em Ann. Univ. Mariae Curie-Sklodowska, Sect. A.} {\bf 35/36} (1982/3)
pp. 19-23.

\bibitem{Barnard5} Barnard, R.W.  ``The omitted area problem for
univalent functions.''  {\em Contemporary Math.} {\bf 38} (1985) pp.
53-60.

\bibitem{Barnard2} Barnard, R.W. ``On bounded univalent functions whose ranges contain
a fixed disk." {\em Trans. Amer. Math. Soc.} {\bf 225} (1977) pp 1-28.

\bibitem{Barnard1} Barnard, R.W.\ ``Open problems and conjectures in
complex analysis.'' in {\em Computational Methods and Function
Theory}, Lecture Notes in Math.\ No.\ 1435 (Springer-Verlag, 1990),
1-26.

\bibitem{Barnard6} Barnard, R.W.  ``A variational technique for bounded
starlike functions.''  {\em Canadian Math. J.}  {\bf 27} No. 2,
(1975) pp. 337-347.

\bibitem{BarnardLewis2} Barnard, R.W., and Lewis, J.L. ``On the omitted area problem."
{\em Mich. Math. J.} {\bf 34} (1987) pp. 13-22.

\bibitem{BarnardLewis1} Barnard, R.W., and Lewis, J.L.  ``Subordination
theorems for some classes of starlike functions.''  {\em Pacific J.
Math.}  {\bf 56} (1975) pp. 333-366.

\bibitem{BarnardPearce1} Barnard, R.W., and Pearce, K.  ``Rounding corners of
gearlike domains and the omitted area problem.''  {\em J.
Comput. Appl. Math.}  {\bf 14} (1986) pp. 217-226.

\bibitem{BarnardPearce2} Barnard, R.W. and Pearce, K. ``Sharp bounds on the
$H_p$ means of the derivative of a convex function for $p=-1$'', {\em
Complex Variables, Theory and Applications} {\bf 21} (1993), 149-158.

\bibitem{BarnardPearce3} Barnard, R.W., and Pearce, K.  ``A short proof of a conjecture
on the integral means of the derivative of a convex function" {\em Ann. Univ. Mariae
Curie-Sklodowska, Sect. A.} {\bf 48} (1994) pp. 1-5.

\bibitem{BarnardSchober1} Barnard, R.W., and Schober, G., ``M\"{o}bius
transformations for convex mappings,'' {\em Complex Variables, Theory
and Applications} {\bf 3} (1984) pp. 45-54.

\bibitem{BarnardSchober2} Barnard, R.W., and Schober, G., ``M\"{o}bius
transformations for convex mappings II,'' {\em Complex Variables,
Theory and Applications} {\bf 7} (1986) pp. 205-214.

\bibitem{BarnardSchober3} Barnard, R.W., and Schober, G., ``M\"{o}bius
transformations of starlike mappings." {\em Complex Variables,
Theory and Applications} {\bf 15} (1990) pp. 211-221.


\bibitem{Bernardi1} Bernardi, S.D.  ``A survey of the development of the
theory of schlicht functions.''  {\em Duke Math. J.}  {\bf 19} (1952) pp.
263-287.

\bibitem{BrannanClunie1} Brannan, D.A., and Clunie, J.G. (eds.).  {\em Aspects
of contemporary complex analysis.}  (Durham, 1979), Academic
Press, London, (1980).

\bibitem{Duren1} Duren, P.  {\em Univalent Functions.}
Spring-Verlag, Vol. 259, New York, (1980).

\bibitem{Goodman1} Goodman, A.W.  ``Note on regions omitted by
univalent functions.''  {\em Bull. Amer. Math. Soc.}  {\bf 55} (1949) pp.
363-369.

\bibitem{Goodman2} Goodman, A.W.  {\em Univalent Functions.}  Mariner
(1982).

\bibitem{GoodmanReich1} Goodman, A.W., and Reich, E.  ``On the regions
omitted by univalent functions II.''  {\em Canad. J. Math.}  {\bf 7}
(1955) pp. 83-88.

\bibitem{Haramard1} Hadamard, J.``Le\c cons sur le calcul des variations."
Paris, 1910. FM pp. 41-432.

\bibitem{Hall1} Hall, R. R. ``On a conjecture of Clunie and
Sheil-Small.''  {\em Bull, London Math. Soc.} {\bf 12} (1980) pp.
25-28.

\bibitem{Hayman1} Hayman, W.K.  {\em Multivalent functions.}  
Cambridge Tracts in Math. and Math. Phys.,  No. 48, Cambridge
University Press, Cambridge, MR21 \#7301 (1958).

\bibitem{Hummel2} Hummel, J. {\em Lectures on variational methods in the
theory of univalent functions.}  No. 8, University of Maryland, 1972.

\bibitem{Hummel1} Hummel, J. ``A variational method for starlike functions."
{\em Proc. Amer. Math. Soc.} {\bf 9} (1958) pp. 82-87.

\bibitem{Jenkins1} Jenkins, J.A.  ``On values omitted by univalent
functions.''  {\em Amer. J. Math.} {\bf 2} (1953) pp. 406-408.

\bibitem{Krzyz1} Krzyz, J. ``Distortion theorems for bounded convex domains II."
{\em Ann. Univ. Mariae Curie-Sklodowska, Sect. A.} {\bf 14} (1960) pp 1028-1030.

\bibitem{Lewandowski1} Lewandowski, Z.  ``On circular symmetrization of
starshaped domains.''  {\em Ann. Univ. Mariae Curie-Sklodowska, Sect.
A.}  {\bf 17} (1963) pp. 35-38.

\bibitem{Lewis1} Lewis, J.L.  ``On the minimum area problem.''  {\em
Indiana Univ. Math. J.}  {\bf 34} (1985) pp. 631-661.

\bibitem{Ma1} Ma, J. ``Julia variations for extremal problems of convex univalent functions."
Ph.D. Dissertation.

\bibitem{MacGregor1} MacGregor, T.H.  ``Geometric problems in complex
analysis.''  {\em Amer. Math. Monthly.}  {\bf 79} (1972) pp.
447-468.

\bibitem{Miller1} Miller, S. (ed.).  {\em Complex analysis.}
(Brockport, NY, 1976), Lecture Notes in Pure and Appl. Math., 
36, Dekker, New York, (1978).

\bibitem{Mocanu1} Mocanu, P.T. ``Une propri\'et\'e de convexit\'e ge\'en\'eralis\'ee
la th\'eorie de la repr\'esentation conforme."  {\em Mathematica (Cluj)} {\bf 11}(34)
(1969) pp. 127-133.

\bibitem{Netanyahu1} Netanyahu, E. ``On univalent functions in the unit disk whose
image contains a given disk."  {\em J. D'analyse Math.} {\bf 23} (1970) pp. 305-322.

\bibitem{Robertson1} Robertson, M.S. ``Extremal problems for analytic functions with positive
real part and applications." {\em Trans. Amer. Math. Soc.} {\bf 106} (1963) pp. 236-253.

\bibitem{SchSchSp1} Schaeffer, A.C., Schiffer, M., and Spencer, D.C. ``The coefficient regions
of schlicht functions." {\em Duke Math J.} {\bf 16} (1949) pp. 493-527.

\bibitem{Schiffer1} Schiffer, M. ``Some recent developments in the theory of
conformal mapping." Appendix in {\em Dirichlet's Principle}, R. Courant.  Interscience , New York (1967).

\bibitem{Stankiewicz1} Stankiewicz, J.  ``On a family of starlike
functions.''  {\em Ann. Univ. Mariae Curie-Sklodowska, Sect. A.}
22-24, (1968-70) pp. 175-181.

\bibitem{Suffridge1} Suffridge, T.  ``A coefficient problem for a class
of univalent functions.''  {\em Mich. Math. J.}  {\bf 16} (1969) pp
33-42.

\bibitem{Valentine1} Valentine, F.A.\ {\em Convex Sets}, McGraw-Hill, New
York, 1964.

\bibitem{Waniurski1} Waniurski, J.  ``On values omitted by convex
univalent mappings.''  {\em Complex Variables, Theory and Appl.}  
{\bf 8} (1987) pp. 173-180.

\bibitem{Warshawski1} Warshawski, S.E. ``On Hadamard's variation formula for 
Green's function."  {\em J. Math. Mech.} {\bf 9} (1960) pp. 497-512.

\bibitem{Zheng1} Zheng, J.\ ``Some extremal problems involving $n$ points
on the unit circle'', Dissertation, Washington University (St.\
Louis), 1991.


\end{thebibliography}


\end{document}
